% ID: SEL-D0023-OA-E04
% ESTADO: validada. S024, cotejada con el examen.
\textbf{a)} Contrastamos $H_0:\mu\le11.7$ frente a $H_1:\mu>11.7$.
La desviación típica es $\sqrt{2.25}=1.5$ cm. Con $n=10$ y
$\bar x=125.8/10=12.58$ cm, el estadístico en la frontera de $H_0$ es
\[
Z=\frac{\bar X-11.7}{1.5/\sqrt{10}},\qquad z_{\mathrm{obs}}\approx1.8552.
\]
Para una media nula menor que 11.7, la probabilidad de caer en
la cola derecha es menor. En la frontera, al $5\%$, se rechaza
si $Z>1.6449$, o equivalentemente si $\bar X>12.4802$ cm
(con umbrales aproximados). La observación cae en la región
de rechazo y aporta evidencia de que la media supera 11.7 cm.

\textbf{b)} Al $3\%$ el cuantil es $1.8808$ y se rechaza si
$\bar X>12.5921$ cm, aproximadamente. Ahora $12.58$ queda fuera
de la región de rechazo. Los datos son compatibles con la
sospecha a este nivel, aunque el contraste no demuestra que sea cierta.
